% Work-MAT | 2. rocnik | hodina 19
\documentclass[11pt,a4paper]{article}
\input{../../sablony/pripravy/priprava.sty}
\usepackage{tikz}
\usetikzlibrary{arrows.meta}
\nastavPripravu{2. ročník}{Exponenciální rovnice}{19}{Exponenciální rovnice IV: kvadratická substituce}
\begin{document}
\poznamka{Navážeme na označování $t=a^x$ z minulé hodiny a využijeme kvadratické rovnice prvního ročníku. Neprohazuj roli $t$ a $x$: nejdříve dokonči pomocnou rovnici, vyluč nepřípustné $t\le0$ a až poté zpětně substituuj. Orientačně 7 min motivace, 15 min ukázka, 13 min procvičení, 10 min kontrola.}
\blok{Když se objeví $a^{2x}$ i $a^x$}
Pravidlo $(a^x)^2=a^{2x}$ nám dovolí označit $t=a^x$. Pak $a^{2x}=t^2$ a exponenciální rovnice se může změnit na kvadratickou. To je výhodné, protože už známe rozklad na součin i diskriminant.
\begin{zadaniBox}[title={Věta - Kvadratická substituce}]
Při substituci $t=a^x$ s $a>0$, $a\ne1$, platí $t>0$ a $a^{2x}=t^2$. Z pomocné kvadratické rovnice je třeba ponechat jen kořeny $t>0$ a pro každý zvlášť dopočítat $x$.
\end{zadaniBox}
\begin{prikladReseniBox}[title={Příklad 1}]
Vyřeš $4^x-6\cdot2^x+8=0$.
\tcblower
\[\begin{aligned}t&=2^x,\quad t>0,\\4^x&=(2^x)^2=t^2,\\t^2-6t+8&=0,\\(t-2)(t-4)&=0.\end{aligned}\]
Oba pomocné kořeny jsou kladné: $t_1=2$, $t_2=4$. Návrat:
\[2^x=2\Rightarrow x_1=1,\qquad 2^x=4\Rightarrow x_2=2.\]
Zkouška: pro $x=1$ je $4-12+8=0$; pro $x=2$ je $16-24+8=0$. $K=\{1,2\}$.
\end{prikladReseniBox}
\begin{prikladReseniBox}[title={Příklad 2}]
Vyřeš $9^x+2\cdot3^x-15=0$.
\tcblower
\[\begin{aligned}t&=3^x>0,\\t^2+2t-15&=0,\\(t+5)(t-3)&=0.\end{aligned}\]
Kořen $t=-5$ vyloučíme, protože $3^x>0$. Pro $t=3$ platí $3^x=3^1$, tedy $x=1$. Zkouška: $9+6-15=0$. $K=\{1\}$.
\end{prikladReseniBox}
\poznamka{Žákům vyzdvihni skutečnost, že množina kořenů pomocné rovnice obecně není množinou kořenů původní rovnice. Zejména $t=-5$ není odpověď na původní otázku.}
\begin{prikladReseniBox}[title={Příklad 3}]
Vyřeš $9^x-4\cdot3^x+3=0$.
\tcblower
$t=3^x>0$: $t^2-4t+3=(t-1)(t-3)=0$. Návrat: $3^x=1=3^0$ nebo $3^x=3=3^1$. Tedy $K=\{0,1\}$.
\end{prikladReseniBox}
\begin{zadaniReseniBox}[title={Samostatná práce 1}]
\begin{enumerate}[label=\alph*)]
\item $4^x-5\cdot2^x+4=0$;
\item $9^x-10\cdot3^x+9=0$;
\item $4^x+2\cdot2^x-8=0$.
\end{enumerate}
\tcblower
\begin{enumerate}[label=\alph*)]
\item $t^2-5t+4=(t-1)(t-4)$, $t=1,4$, tedy $x=0,2$.
\item $t^2-10t+9=(t-1)(t-9)$, $t=1,9$, tedy $x=0,2$.
\item $t^2+2t-8=(t+4)(t-2)$, záporné $t=-4$ vyloučíme, z $t=2$ je $x=1$.
\end{enumerate}
\end{zadaniReseniBox}
\begin{zadaniReseniBox}[title={Rozšíření 1}]
Vyřeš $16^x-5\cdot4^x+4=0$.
\tcblower
$t=4^x>0$. Potom $t^2-5t+4=(t-1)(t-4)=0$. Z $4^x=1$ a $4^x=4$ plyne $K=\{0,1\}$.
\end{zadaniReseniBox}
\blok{Výběr metody před dalším procvičováním}
Než začneš počítat, vždy rozhodni, zda jde obě strany rovnice převést na \emph{společný základ}, zda lze \emph{vytknout} stejný exponenciální výraz, nebo zda se opakuje dvojice $a^x$ a $a^{2x}$ vhodná pro \emph{kvadratickou substituci}. V další části učiva budeme tyto postupy kombinovat.
\begin{zadaniReseniBox}[title={Samostatná práce 2}]
U každé rovnice nejprve vyber nejvhodnější metodu a poté ji vyřeš: a) $8^x=4^{x+1}$; b) $2^{x+1}+2^x=12$; c) $4^x-5\cdot2^x+4=0$.
\tcblower
\begin{enumerate}[label=\alph*)]
\item Společný základ $2$: $3x=2x+2$, takže $x=2$.
\item Vytýkání: $3\cdot2^x=12$, takže $x=2$.
\item Substituce $t=2^x>0$: $(t-1)(t-4)=0$, takže $x=0$ nebo $x=2$.
\end{enumerate}
\end{zadaniReseniBox}
\begin{zadaniReseniBox}[title={Domácí úkol}]
\begin{enumerate}[label=\alph*)]
\item Vyřeš $4^x-3\cdot2^x+2=0$.
\item Vyřeš $9^x+4\cdot3^x-5=0$.
\item Proč není $t=-1$ přípustné řešení substituce $t=3^x$?
\end{enumerate}
\tcblower
\begin{enumerate}[label=\alph*)]
\item $(t-1)(t-2)=0$, tedy $K=\{0,1\}$.
\item $(t+5)(t-1)=0$; jen $t=1>0$, takže $K=\{0\}$.
\item Protože $3^x>0$ pro všechna reálná $x$.
\end{enumerate}
\end{zadaniReseniBox}
\end{document}
